% =====================================================================
% MDBC --- Supplementary Information
% Build: pdflatex -> bibtex -> pdflatex x2  (shares references.bib)
% =====================================================================
\documentclass[11pt]{article}

\usepackage[utf8]{inputenc}
\usepackage[T1]{fontenc}
\usepackage{lmodern}
\usepackage[a4paper,margin=1in]{geometry}
\usepackage{graphicx}
\usepackage{amsmath,amssymb}
\usepackage{booktabs}
\usepackage{longtable}
\usepackage{siunitx}
\usepackage{xcolor}
\usepackage[hidelinks]{hyperref}
\usepackage[round,authoryear]{natbib}

\newcommand{\TODO}[1]{\textcolor{red}{[\textbf{TODO:} #1]}}
\newcommand{\PLACE}[1]{\textcolor{blue}{[\textit{#1}]}}

% number SI floats as S1, S2, ...
\renewcommand{\thetable}{S\arabic{table}}
\renewcommand{\thefigure}{S\arabic{figure}}
\renewcommand{\theequation}{S\arabic{equation}}
\renewcommand{\thesection}{S\arabic{section}}

\title{\textbf{Supplementary Information}\\[4pt]
\large Cross-delta transferability limits of mangrove blue-carbon stock models}
\author{Md Rakib Hasan}
\date{\today}

\begin{document}
\maketitle
\tableofcontents
\bigskip

% ---------------------------------------------------------------------
\section{Delta selection and the frozen benchmark}
\PLACE{Full pre-registered criteria (C1--C5), candidate windows, the GMW-clip
procedure, the usable-SOC-core gate, and the change log (estuarine systems added
for continental coverage; the Oceania data-gap decision). Mirror
\texttt{docs/stage0\_delta\_selection\_criteria.md}.}

\begin{longtable}{llrrrl}
\caption{Full candidate delta registry: mangrove area (GMW v3, 2020, equal-area),
usable SOC cores, raw in-window cores, and assigned role.
\PLACE{regenerate from \texttt{data/processed/delta\_registry.csv}.}}\\
\toprule
Delta & Country & Area (km$^2$) & SOC cores & Raw cores & Role \\
\midrule
\endfirsthead
\toprule
Delta & Country & Area (km$^2$) & SOC cores & Raw cores & Role \\
\midrule
\endhead
\PLACE{populate all 16 windows} & & & & & \\
\bottomrule
\end{longtable}

% ---------------------------------------------------------------------
\section{SOC label harmonization}
\PLACE{Detail Eq.~\ref{eq:soc-si}: layer cleaning bounds (BD, $f_C$ ranges),
OM$\to$C ratio derivation ($\approx0.427$ over $\sim$12k paired CCN samples),
$\geq$80~cm coverage requirement, deepest-layer extrapolation to 100~cm and its
flag, 0.5/99.5 percentile outlier trim. Report final $n$, median and range.}
\begin{equation}
\mathrm{SOC}_{0\text{--}100} = 100 \sum_i \mathrm{BD}_i\, f_{C,i}\, \Delta z_i,
\qquad f_{C,i} = \begin{cases} f_{C,i}^{\text{meas}} \\ 0.427\, f_{OM,i} \end{cases}
\label{eq:soc-si}
\end{equation}

% ---------------------------------------------------------------------
\section{Covariate stack}
\PLACE{Full covariate list with source, native resolution, CRS and sampling
method (nearest finite cell). Note the PEDOFLUX-lake layers that were truncated
and excluded (SoilGrids texture, OpenLandMap, MERIT terrain) and why GSOCmap is
treated cautiously (circular prior). EOT20 constituents (M2, S2, K1, O1) and the
tidal-range / form-factor definitions.}

\begin{table}[h]
\centering
\caption{Covariates used in the v1 SOC model and per-core coverage.
\PLACE{from the \texttt{sample\_covariates.py} coverage report.}}
\begin{tabular}{lllr}
\toprule
Covariate group & Source & Variables & Coverage \\
\midrule
Climate    & CHELSA v2.1   & bio1--19            & \PLACE{100\%} \\
SOC prior  & GSOCmap        & stock               & \PLACE{77\%} \\
Soil       & SoilGrids      & bdod, cec (0--5 cm) & \PLACE{58\%} \\
Land cover & Copernicus LULC& class (2015)        & \PLACE{100\%} \\
Geomorphic & Natural Earth  & dist.\ coast/river  & \PLACE{100\%} \\
Tidal      & EOT20          & range mean/spring, form & \PLACE{100\%} \\
\bottomrule
\end{tabular}
\end{table}

% ---------------------------------------------------------------------
\section{Validation, AOA and conformal details}
\PLACE{Spatial-block size relative to the SOC autocorrelation range; LODO fold
definition and leakage control; AOA dissimilarity-index formula and threshold
choice; split-conformal construction and the $\alpha=0.1$ target. Sensitivity of
the AOA threshold.}

% ---------------------------------------------------------------------
\section{Full per-delta transfer results}
\PLACE{Per-delta LODO $R^2$, RMSE, bias, AOA-inside fraction and conformal
coverage for both models and all three feature sets (climate; +geomorphic;
+tidal). Mirror \texttt{soc\_lodo\_results.json}.}

\begin{table}[h]
\centering
\caption{Per-delta leave-one-delta-out results (gradient boosting, full feature
set). \PLACE{from \texttt{soc\_lodo\_results.json}.}}
\begin{tabular}{lrrrrr}
\toprule
Delta & $n$ & $R^2$ & RMSE (Mg ha$^{-1}$) & AOA in & Conf.\ cov. \\
\midrule
\PLACE{populate 8 core deltas} & & & & & \\
\bottomrule
\end{tabular}
\end{table}

% ---------------------------------------------------------------------
\section{Transfer diagnostic}
\PLACE{Energy-distance and dissimilarity-index methodology; the full per-delta
shift table; Spearman correlations with the small-$n$ power caveat; the
standardized-mean-difference ranking that names each delta's dominant shift axis.
Mirror \texttt{docs/stage2\_diagnostic.md}.}

% ---------------------------------------------------------------------
\section{Aboveground and total ecosystem carbon (reserved)}
\TODO{Simard 2019 AGB acquired for all core deltas. To add: AGB$\to$C labels,
AGB LODO models, AGB-C + SOC total with uncertainty propagation, and per-delta
total-carbon maps with AOA masks.}

% ---------------------------------------------------------------------
\section{Reproducibility}
\PLACE{Pinned environment (\texttt{requirements.txt}), single-command runner
(\texttt{run\_all.sh}), the \texttt{MDBC\_LAKE} external-driver convention, and
the provenance manifests. Frozen LODO splits released with the benchmark.}

\bibliographystyle{plainnat}
\bibliography{references}

\end{document}
